\begin{afterword}
\summaryhead

The post asks what physicists thought before decoherence, and reconstructs the first idea of collapse: squared amplitudes were read as probabilities, and learning that something did not happen was taken to make its probability zero. Read literally, the author says, this makes knowledge or consciousness the cause of collapse, which was Heisenberg's theory. Later theories made collapse an objective process that keeps one part of the wavefunction and acts before superpositions reach human size.

The author asks why collapse theories keep a single survivor, and answers that they were devised when no one imagined more than one world. Experiments, the post says, keep finding superposition in larger systems, so why does no one propose collapse at planetary scale? A cynic's answer, which the post then calls ad hominem, is that collapse preserves the intuition of a single self. No collapse theory has experimental support. The post ends with eight ways in which collapse would be unique in physics, among them non-linear, non-unitary, discontinuous and faster than light.

\responsehead

The eight-item list is, for the most part, an accurate description of objective collapse. Collapse theories do add nonlinear and stochastic terms to the Schrödinger equation, and their proponents say so. The complaint behind the list, that the textbook postulate is vague, is shared by those proponents, who call it ``absolutely unacceptable.'' The problems are in the post's history and in what it leaves out of the comparison.

The history is loosely told in three places. The post says it had not occurred to anyone that the quantum laws might apply universally. In fact applying them to the measuring apparatus is what produced the measurement problem, and collapse was added to avoid the result. The post attributes to Heisenberg a theory in which knowledge ``or even conscious awareness'' causes collapse; Heisenberg tied the probability function to knowledge, but wrote that the transition to the actual ``is not connected with the act of registration of the result by the mind of the observer.'' And the post says collapse theories were devised when no physicist imagined more than one world. That is true of the 1932 postulate, not of the theories ``considered in modern academia,'' which date from 1976 and 1986, after Everett.

Some items in the list are weaker than they look. The discontinuity of item 3 does not hold for the continuous collapse model. The CPT violation of item 5 is asserted, and a later argument holds that collapse models can be time-symmetric. The nonlocality of item 8 is shared by every theory in which measurements have single outcomes, and collapse sends no signals; it is ``the only'' nonlocal phenomenon only if many-worlds, or some other way around Bell's theorem, is already assumed. The experimental record also cuts less than the post says. No experiment supports collapse, but none separates many-worlds from standard quantum mechanics either, while collapse theories make predictions that can be tested and are being tested.

The question about collapse at planetary scale has the answer the post itself gave: collapse theories are meant to give measurements single outcomes, and a theory with superposed observers does not. The post instead offers a cynic's reading of motives, and itself calls it ad hominem.

In short: an accurate list of what objective collapse gives up, set in a loosely told history, with the costs of the rival view left for a later post.
\end{afterword}
